\documentclass[11pt]{article}
\usepackage[localtheorems]{raeez-math-template}


\theoremstyle{plain}
\newtheorem{theorem}{Theorem}
\newtheorem{proposition}[theorem]{Proposition}
\newtheorem{lemma}[theorem]{Lemma}
\newtheorem{corollary}[theorem]{Corollary}

\theoremstyle{definition}
\newtheorem{definition}[theorem]{Definition}
\newtheorem{construction}[theorem]{Construction}

\theoremstyle{remark}
\newtheorem{remark}[theorem]{Remark}

\newcommand{\C}{\mathbb{C}}
\newcommand{\R}{\mathbb{R}}
\newcommand{\Z}{\mathbb{Z}}
\newcommand{\N}{\mathbb{N}}
\newcommand{\Q}{\mathbb{Q}}
\newcommand{\F}{\mathbb{F}}
\newcommand{\HH}{\mathrm{HH}}
\newcommand{\Rep}{\mathrm{Rep}}
\newcommand{\Vect}{\mathrm{Vect}}
\newcommand{\Hom}{\mathrm{Hom}}
\newcommand{\End}{\mathrm{End}}
\newcommand{\id}{\mathrm{id}}
\newcommand{\tr}{\mathrm{tr}}
\newcommand{\Tot}{\mathrm{Tot}}
\newcommand{\Sym}{\mathrm{Sym}}
\newcommand{\fg}{\mathfrak{g}}
\newcommand{\fh}{\mathfrak{h}}
\newcommand{\fsl}{\mathfrak{sl}}
\newcommand{\fgl}{\mathfrak{gl}}
\newcommand{\fL}{\mathfrak{L}}
\newcommand{\cA}{\mathcal{A}}
\newcommand{\cB}{\mathcal{B}}
\newcommand{\cC}{\mathcal{C}}
\newcommand{\cE}{\mathcal{E}}
\newcommand{\cF}{\mathcal{F}}
\newcommand{\cM}{\mathcal{M}}
\newcommand{\cO}{\mathcal{O}}
\newcommand{\CE}{\mathrm{CE}}
\newcommand{\Conf}{\mathrm{Conf}}
\newcommand{\kch}{\kappa_{\mathrm{ch}}}
\newcommand{\kcat}{\kappa_{\mathrm{cat}}}
\newcommand{\kBKM}{\kappa_{\mathrm{BKM}}}
\newcommand{\kfib}{\kappa_{\mathrm{fiber}}}
\newcommand{\Ucha}{U^{\mathrm{ch}}}

\title{Class $\mathcal{M}$: Reduced $E_3$ Bar Cohomology Through Genus Three}
\author{Raeez Lorgat}
\date{22 April 2026}

\begin{document}

\maketitle

\begin{abstract}
Let $A$ be a chiral algebra of shadow class $\mathcal{M}$ whose reduced
external $E_3$ one-handle bar page is
$\Lambda^\bullet(\C^3)$. The reduced genus-$g$ handle-copy chain page has
Poincar\'e series
$((1+s)(1+t)(1+u))^g$ and dimension $8^g$. If the quartic shadow component
acts by the rank-one pairing between exterior degrees $0$ and $3$, with
nonzero coefficient $8\kch(A)S_4(A)$, the post-$d_4$ page has Poincar\'e
series $(3t+3t^2)^g$ and dimension $6^g$. Degree considerations prove
$E_4=E_\infty$ for $g\leq 3$, giving cohomological dimensions
$6,36,216$. For $g\geq 4$, $6^g$ is only the post-$d_4$ dimension; the
first uncomputed differential is $d_5$, and the final cohomology dimension
requires its rank and the ranks of the later admissible differentials.
\end{abstract}

\tableofcontents

\section{The Class \texorpdfstring{$\mathcal{M}$}{M} Condition}
\label{sec:what-M-names}

The symbol $\mathcal{M}$ denotes the fourth class in the shadow-tower
quadrichotomy
$\mathcal{G}/\mathcal{L}/\mathcal{C}/\mathcal{M}$. A chiral algebra $A$ is
class $\mathcal{M}$ (\emph{mixed}, or \emph{mock} in the moonshine
specialisation) when its shadow invariants $S_k$, defined by the
Maurer--Cartan recursion on its $A_\infty$-coproduct
\[
 m(a) \;=\; \sum_{k\ge 2} S_k \cdot \Theta_k(a),
\]
satisfy $S_k\ne 0$ for \emph{all} $k\ge 3$. Class $\mathcal{G}$ is Gaussian
($S_k=0$ for $k\ge 3$, Heisenberg); class $\mathcal{L}$ is lemniscatic
($S_3\ne 0$, $S_{k\ge 4}=0$, Yangian); class $\mathcal{C}$ has finitely many
nonzero $S_k$ ($bc$-system). The archetype of $\mathcal{M}$ is Virasoro at
$c=1$, where $S_4=10/27$, $S_3=1/2$, and the sign-alternating tower never
terminates.

The four readings of $\mathcal{M}$ are compatible with this single definition:
mixed Gevrey-$1$ shadow series, mock-modular twining genera, Mathieu
$M_{24}$ symmetry in the K3 elliptic genus, and moonshine denominators at
central charge $24$. The defining condition is the nonvanishing shadow
tower. The moonshine and K3 interpretations supply examples and
representations of the same tower.

The K3 sigma-model vertex algebra $\cA_\sigma$ at a generic moduli point is
the governing example. With $c=6$, $\cN=(4,4)$ superconformal symmetry, and
shadow data $(\kch, S_3, S_4)=(2,\,2,\,5/156)$, the discriminant
$\Delta=8\kch S_4=20/39\ne 0$ forces $S_r\ne 0$ for every $r\ge 5$ by the
shadow recurrence. Hence $\cA_\sigma$ is class $\mathcal{M}$ and
$\kch(\cA_\sigma)=2=\dim_\C K3$.

\section{The Reduced External \texorpdfstring{$E_3$}{E3} Bar Tricomplex}
\label{sec:e3-bar-chain-level}

For a CY$_d$ input with $d\geq 3$, the CY-to-chiral functor produces an
$E_1$-chiral algebra $A$. The $E_2$ structure lives on the centre or module
layer, for example on $Z^{\mathrm{der}}_{\mathrm{ch}}(A)$ and on
$\Rep(A)$. The $E_3$ in this note is external: it is the three-direction
Hochschild/factorisation tricomplex coming from configuration-space
operations on $\C^3$.

Five objects remain separate throughout:
\[
 A,\qquad B(A),\qquad A^i,\qquad A^!,\qquad
 Z^{\mathrm{der}}_{\mathrm{ch}}(A).
\]
The identity $\Omega B(A)\simeq A$ is bar--cobar inversion. The object
$A^!$ is the Verdier/Koszul dual. The bulk algebra is computed by
Hochschild cochains, equivalently by the derived chiral centre.

\begin{definition}[The $E_3$ bar tricomplex]
\label{def:e3-tricomplex}
Let $V\simeq\C$ be the reduced one-handle augmentation quotient of the
external three-direction bar construction. The one-handle $E_3$ bar
tricomplex is
\[
 B_{E_3}^{(1)}(A)^{p,q,r} \;=\; \Lambda^p(V)\otimes\Lambda^q(V)\otimes\Lambda^r(V),
\]
with three commuting differentials whose Omega-background weights are the
equivariant parameters $(h_1,h_2,h_3)$, subject to the Calabi--Yau constraint
$h_1+h_2+h_3=0$.
The reduced genus-$g$ model is the tensor power
\[
 B_{E_3}^{(g)}(A) \;=\; \bigl(B_{E_3}^{(1)}(A)\bigr)^{\otimes g}.
\]
\end{definition}

For one handle in the reduced one-generator model, the tridegrees are the
eight exterior monomials of $\Lambda^\bullet(\C^3)$:
\[
 (0,0,0),\ (1,0,0),\ (0,1,0),\ (0,0,1),\
 (1,1,0),\ (1,0,1),\ (0,1,1),\ (1,1,1).
\]
For the reduced genus-$g$ model, tensor multiplication gives
\[
 P_{\mathrm{chain}}(s,t,u) \;=\; \bigl((1+s)(1+t)(1+u)\bigr)^g,
 \qquad P_{\mathrm{chain}}(1,1,1) = 8^g.
\]
At $g=1,2,3$ the chain dimensions are $8,64,512$.

\section{The Quartic Differential}
\label{sec:d4-and-6g}

For class $\mathcal{G}$, all differentials vanish and the cohomology equals
the chain page. For class $\mathcal{L}$, the shadow tower terminates before
the quartic term. For class $\mathcal{C}$, charge conservation kills the
quartic contribution. Class $\mathcal{M}$ has $S_4\ne 0$, so the quartic
shadow supplies the first differential beyond the three coordinate
differentials.

\begin{proposition}[Virasoro $d_4$, per-factor cohomology $3t+3t^2$]
\label{prop:virasoro-d4-per-factor}
For the Virasoro chiral algebra at central charge $c=1$, the $E_4$-page of
the one-handle tricomplex has Poincar\'e polynomial
\[
 P_{E_4}^{\mathcal{M},\,n=1}(t) \;=\; 3t + 3t^2,
 \qquad \dim P_{E_4}^{\mathcal{M},\,n=1} = 6.
\]
The differential $d_4$ has weight
$\Delta=8\kch S_4=8\cdot\tfrac{1}{2}\cdot\tfrac{10}{27}
=\tfrac{40}{27}$.
\end{proposition}

\begin{proof}
The chain-level tridegree polynomial for $n=1$ is
$(1+s)(1+t)(1+u)$. The quartic shadow gives a rank-one pairing between
the exterior-degree $0$ and exterior-degree $3$ classes. After setting
$s=t=u$ on the post-$d_4$ page, the surviving classes are
\[
 \{(1,0,0),(0,1,0),(0,0,1),(1,1,0),(1,0,1),(0,1,1)\},
\]
three in total degree $1$ and three in total degree $2$. Thus
$P_{E_4}(t)=3t+3t^2$. The coefficient
$\Delta=8\kch S_4$ is nonzero at Virasoro $c=1$ because
$\kch=1/2$ and $S_4=10/27$, giving $\Delta=40/27$.
\end{proof}

\begin{theorem}[Class $\mathcal{M}$ $E_3$ bar cohomology through genus three]
\label{thm:class-m-e3-bar-6g}
Let $A$ be a class-$\mathcal{M}$ chiral algebra whose reduced external $E_3$
one-handle bar page is the reduced model above. Assume that its quartic
shadow component acts as the rank-one pairing
\[
 \Lambda^3(\C^3)\longrightarrow\Lambda^0(\C^3),
 \qquad \Delta(A)=8\kch(A)S_4(A)\ne 0.
\]
For $1\leq g\leq 3$,
\[
 E_4\bigl(B_{E_3}^{(g)}(A)\bigr)=
 E_\infty\bigl(B_{E_3}^{(g)}(A)\bigr),
 \qquad
 P_{\mathrm{coh}}^{\mathcal{M}}(t;g) \;=\; \bigl(3t+3t^2\bigr)^g,
\]
with total dimension
\[
 \dim H^\bullet\bigl(B_{E_3}^{(g)}(A)\bigr) \;=\; 6^g.
\]
The finite chain-to-cohomology deficit is $8^g-6^g$.
\end{theorem}

\begin{proof}
The reduced genus-$g$ model is the tensor product of $g$ copies of the
one-handle post-$d_4$ complex. The rank-one quartic hypothesis gives the
same per-handle survivor list as Proposition
\ref{prop:virasoro-d4-per-factor}. Hence the per-handle polynomial is
$3t+3t^2$, and K\"unneth gives $(3t+3t^2)^g$ with total dimension $6^g$ on
the post-$d_4$ page.

This page is supported in total degrees $[g,2g]$. A differential $d_r$ for
$r\geq 5$ shifts total degree by $r-1\geq 4$. For $g=1,2,3$, there is no
source and target pair inside $[g,2g]$ with shift $4$ or larger. Thus
$d_r=0$ for every $r\geq 5$ in these genera, and $E_4=E_\infty$.
\end{proof}

\begin{center}
\renewcommand{\arraystretch}{1.2}
\begin{tabular}{lcccc}
\toprule
 Shadow class & $g=1$ & $g=2$ & $g=3$ & Formula \\
 \midrule
 $\mathcal{G}$ & $8$ & $64$ & $512$ & $8^g$ \\
 $\mathcal{L},\mathcal{C}$ & $8$ & $64$ & $512$ & $2^{3g}=8^g$ \\
 $\mathcal{M}$ & $6$ & $36$ & $216$ & $6^g$ \\
 Deficit & $2$ & $28$ & $296$ & $8^g-6^g$ \\
 \bottomrule
\end{tabular}
\end{center}

\section{The First Higher-Genus Obstruction}
\label{sec:higher-genus-obstruction}

For $g\geq 4$, under the same reduced-model hypotheses, the proved value
$6^g$ is the post-$d_4$ dimension. The final cohomology dimension requires
the higher differentials. The first obstruction occurs at $g=4$: the
total-degree support is $[4,8]$, and a $d_5$ differential has shift $4$, so
the first map to compute is
\[
 d_5 \colon E_5^{8}\longrightarrow E_5^{4}.
\]
Both source and target lie inside the support. A proof of the genus-$4$
cohomology dimension must compute this map on
\[
 (\Lambda^1\C^3\oplus\Lambda^2\C^3)^{\otimes 4},
\]
including its coefficient from the quintic shadow invariant $S_5$, its
compatibility with the already imposed $d_4$ relation, and its rank.

For general $g\geq 4$, the proof obligation is the finite list
\[
 d_r \colon E_r^k\longrightarrow E_r^{k-r+1},
 \qquad 5\leq r\leq g+1,\qquad g\leq k-r+1<k\leq 2g.
\]
The infinite shadow tower supplies the only possible coefficients for these maps.
In each fixed genus it gives a finite spectral-sequence problem. The cohomology
through genus three is $6^g$; beyond that range the post-$d_4$ value
$6^g$ is the input to the higher-differential computation.

\section{The Chiral CE Identity}
\label{sec:chiral-ce}

The bar complex of the chiral envelope identifies with the
Chevalley--Eilenberg chain complex of the underlying Lie conformal
algebra:

\begin{proposition}[Bar--CE identification]
\label{prop:bar-ce}
For a Lie conformal algebra $\fL$,
\[
 B\bigl(\Ucha(\fL)\bigr) \;\simeq\; \CE_\bullet(\fL)
 \;=\;\bigl(\Lambda^\bullet\fL,\, d_{\CE}\bigr),
\]
with $d_{\CE}$ encoding the $\lambda$-bracket. For an $L_\infty$-conformal
algebra, the identification extends to
\[
 B\bigl(\Ucha_{\mathrm{der}}(\fL)\bigr) \;\simeq\; \CE^{L_\infty}_\bullet(\fL),
 \qquad d_{\CE}^{L_\infty} = \sum_{k\ge 2} d^{(k)},
\]
where $d^{(k)}$ is the differential coming from the structure map $l_k$.
\end{proposition}

On the class-$\mathcal{M}$ side, the $L_\infty$-structure maps $l_k$ for
$k\ge 3$ are nonvanishing, in exact correspondence with the shadow tower
$S_k\ne 0$. The chiral CE differential $d_{\CE}^{L_\infty}$ thus has
components of all weights $k\ge 2$. The first component relevant to the
external $E_3$ tricomplex is the quartic-shadow contribution, realised as
the $d_4$ differential above. In the one-generator Virasoro example the
ordinary exterior algebra cannot carry the whole higher operation by
antisymmetry. The nontrivial $L_\infty$ content lives on the tensor coalgebra
$T^c(s^{-1}T)$, where the shadow tower packages into the tricomplex.

\section{6d Holomorphic Chern--Simons Reading}
\label{sec:6d-hcs-interweave}

Consider 6d holomorphic Chern--Simons theory on a holomorphic threefold
$X_g$ whose boundary reduction has a K3 sector and a genus-$g$ handle
factor. The product $K3\times\Sigma_g$ is a compact Calabi--Yau threefold
only when $\Sigma_g$ is elliptic; for $g\ne 1$ the hCS background is the
local or canonical-twisted threefold carrying the same handle
factorisation. Its boundary observables form a factorisation
algebra with external $E_3$ operations along the three holomorphic
directions. Compactification on $\Sigma_g$ produces factorisation homology
$\int_{\Sigma_g}A$ of the boundary algebra.

\begin{proposition}[Handlebody comparison]
\label{prop:6d-hcs-skein}
Let the boundary algebra of the holomorphic Chern--Simons background
$X_g$ be the class-$\mathcal{M}$ algebra $A$ above, and assume the
handlebody factorisation homology is computed by the reduced external $E_3$ bar
tricomplex. Then for $1\leq g\leq 3$,
\[
 \int_{\Sigma_g} A \;\simeq\; H^\bullet\bigl(B_{E_3}^{(g)}(A)\bigr),
 \qquad \dim_\C\!\int_{\Sigma_g} A \;=\; 6^g.
\]
\end{proposition}

\begin{proof}
The stated hypotheses identify the factorisation homology with the
handlewise $E_3$ bar construction. Theorem
\ref{thm:class-m-e3-bar-6g} computes that construction through genus three.
\end{proof}

\section{Invariant Separations}
\label{sec:cross-consistency}

\begin{itemize}
 \item $\kcat(K3\times E)=0$ by K\"unneth-multiplicativity on Euler
   characteristics: $\chi(\cO_{K3})\cdot\chi(\cO_E)=2\cdot 0=0$.
   The compact chiral value is $\kch(K3\times E)=0$.
   The Heisenberg specialisation has $\kch^{\mathrm{Heis}}=3$.
   The Borcherds denominator has $\kBKM(\mathfrak{g}_{\Delta_5})=5$
   for the $\Delta_5$ form. The K3 fibre lattice has $\kfib=24$.
 \item $\kBKM(\Phi_N)=c_N(0)/2$ for $N\in\{1,2,3,4,6\}$ is the Borcherds
   weight theorem for the CHL BKM denominator forms. The primitive CHL
   constants are
   \[
    (c_1(0),c_2(0),c_3(0),c_4(0),c_6(0))=(10,8,6,4,2),
   \]
   hence
   \[
    \bigl(\kBKM(\Phi_1),\kBKM(\Phi_2),\kBKM(\Phi_3),
    \kBKM(\Phi_4),\kBKM(\Phi_6)\bigr)=(5,4,3,2,1).
   \]
   These automorphic weights are independent of the $E_3$ bar
   cohomological dimension.
 \item The Gritsenko--Cl\'ery diagonal-divisor catalogue is a separate
   eight-row classification. Its twice-weight tuple
   $(10,4,6,2,4,1,3,2)$ is not the primitive CHL tuple
   $(10,8,6,4,2)$, although both catalogues contain the $\Delta_5$ row.
 \item $\mathrm{CoHA}(\C^3)=Y^+$, the positive half. Drinfeld-double and
   Koszul-dual constructions form separate representation-theoretic layers.
 \item The constructions of $G(K3\times E)$ are distinct: Mukai lattice,
   Igusa $\Phi_{10}$ via Gritsenko $\Delta_5$, BKM denominator, K3 fibre
   rank, quantum toroidal algebra, and CoHA. The external $E_3$ bar
   tricomplex is attached to $\cA_\sigma$ and is independent of this list.
\end{itemize}

\section{Open Problems}
\label{sec:open}

\begin{enumerate}
 \item Genus $4$: compute
   $d_5:E_5^8\to E_5^4$ on
   $(\Lambda^1\C^3\oplus\Lambda^2\C^3)^{\otimes4}$, including the
   $S_5$ coefficient and the rank of the induced map.
 \item Higher genus: compute all admissible $d_r$ maps with
   $5\leq r\leq g+1$ and $g\leq k-r+1<k\leq 2g$.
 \item $M_{24}$-equivariant refinement: at $g=1$, the $E_4$ page $6=3\cdot 2$
   splits into a $3$-dimensional degree-$1$ summand and a $3$-dimensional
   degree-$2$ summand. The proof obligation is to construct twined
   insertions on these summands and compare their characters, conjugacy
   class by conjugacy class, with the Mathieu-twined Euler characteristics
   $Z(K3;\tau,z)_{[m]}$, mock Jacobi forms of weight $0$ and index $1$.
 \item Skein realisation: construct a reduced skein-state model of the
   genus-$g$ handlebody whose state space is functorially identified with
   $H^\bullet(B_{E_3}^{(g)}(\cA_\sigma))$ through genus three. The
   dictionary between shadow classes and Kauffman parameters is part of the
   construction.
 \item Mathieu-twined family: construct insertions
   $\{\cA_\sigma^{[m]}\}_{[m]\subset M_{24}}$ on the K3 sigma-model sector
   and prove that the genus-three-and-lower reduced external $E_3$ bar
   cohomology retains the universal class-$\mathcal{M}$ page
   $(3t+3t^2)^g$.
\end{enumerate}

\end{document}
